% ==============================================================
\section{Problem Definition}\label{sec:problem}
% ==============================================================

This section reviews the one-dimensional setting
of Apriori-Window~\cite{aprioriwindow} (\S\ref{sec:1d_review})
and then formulates its $d$-dimensional extension (\S\ref{sec:nd_extension}).

% ==============================================================
\subsection{Apriori-Window (One-Dimensional)}\label{sec:1d_review}
% ==============================================================

Let $\mathcal{I} = \{i_1, \ldots, i_m\}$ be the item universe and
$\mathcal{D}_{\mathrm{1D}} = \{(t_n, d_n)\}_{n=1}^N$
be a timestamped transaction sequence,
where $t_n \in \mathbb{Z}_{\geq 0}$ is an integer timestamp
and $d_n \subseteq \mathcal{I}$ is an itemset.

The occurrence sequence of a pattern $P \subseteq \mathcal{I}$ is
\[
  \mathrm{occ}(P) = \{t_n : P \subseteq d_n\}.
\]

\begin{definition}[Interval support~\cite{aprioriwindow}]\label{def:sup1d}
Under window width $W$ ($W > 0$),
the interval support of pattern $P$
at window left endpoint $l$ is defined as
\[
  \mathrm{sup}(P,\, l,\, W)
  = \bigl|\{t \in \mathrm{occ}(P) : l \leq t \leq l + W\}\bigr|.
\]
The window $[l, l+W]$ contains $W+1$ discrete time points.
\end{definition}

\begin{definition}[Dense-start set and dense interval]\label{def:dense1d}
For a support threshold $\theta$,
\[
  \mathcal{T}_\theta(P)
  = \{l \in \mathbb{Z}_{\geq 0} : \mathrm{sup}(P, l, W) \geq \theta\}
\]
is the \emph{dense-start set} of $P$.

Let $[a,b] \subseteq \mathcal{T}_\theta(P)$
be a maximal contiguous component.
Then $[a,\, b + W]$ is called a \emph{dense interval} of $P$.
\end{definition}


% ==============================================================
\subsection{Extension to $d$ Dimensions}\label{sec:nd_extension}
% ==============================================================

Let $\mathcal{D} = \{(\bm{b}_n, d_n)\}_{n=1}^N$
be a $d$-dimensional transaction sequence, where
\[
  \bm{b}_n = (b_n^{(0)}, b_n^{(1)}, \ldots, b_n^{(d-1)})
  \in \prod_{k=0}^{d-1}\{0,1,\ldots,N_k-1\}
\]
is an integer coordinate vector (axis 0 is time) and
$d_n \subseteq \mathcal{I}$ is an itemset.
$N$ denotes the \emph{number of transactions}.

\begin{definition}[Temporal multiset reduction to a fixed spatial start point]\label{def:collapsed_multiset_nd}
Fix a spatial start vector
\[
  \bm{s}_{\mathrm{sp}} = (s_1, \ldots, s_{d-1})
\]
and spatial window widths
$\bm{W}_{\mathrm{sp}} = (W_1, \ldots, W_{d-1})$
(each $W_k > 0$).
The spatial window along axis $k$, $[s_k, s_k + W_k]$,
contains $W_k + 1$ discrete coordinates.
The temporal multiset of pattern $P$ is defined as
\[
  L_P(\bm{s}_{\mathrm{sp}})
  =
  \{\!\{
    b_n^{(0)}
    \mid
    P \subseteq d_n,\;
    s_k \leq b_n^{(k)} \leq s_k + W_k\; (\forall k=1,\ldots,d-1)
  \}\!\}
\]
where $\{\!\{\cdot\}\!\}$ denotes a multiset allowing duplicates.
\end{definition}

The essence of this definition is that,
for each fixed spatial start point,
the $d$-dimensional problem is reduced to
a \emph{one-dimensional problem with duplicate timestamps}.

\begin{definition}[$d$-dimensional sliding window support]\label{def:support_nd}
For a temporal start point $s_0$ and spatial start vector
$\bm{s}_{\mathrm{sp}}$,
the support of pattern $P$ is defined as
\begin{equation}\label{eq:support_nd}
  S_P(s_0, \bm{s}_{\mathrm{sp}})
  = \bigl|\{t \in L_P(\bm{s}_{\mathrm{sp}}) : s_0 \leq t \leq s_0 + W_0\}\bigr|.
\end{equation}
\end{definition}

\begin{definition}[Dense-start set and dense region]\label{def:dense_region_nd}
For a fixed spatial start vector $\bm{s}_{\mathrm{sp}}$,
\[
  \mathcal{T}_\theta(P,\bm{s}_{\mathrm{sp}})
  =
  \{s_0 \in \mathbb{Z}_{\geq 0} : S_P(s_0,\bm{s}_{\mathrm{sp}}) \geq \theta\}
\]
is the \emph{dense-start set} of $P$
at $\bm{s}_{\mathrm{sp}}$.

Let $[a_j,b_j]$ be a maximal contiguous component
of $\mathcal{T}_\theta(P,\bm{s}_{\mathrm{sp}})$.
Then $[a_j,\, b_j + W_0]$ is called a \emph{dense interval} of $P$
at $\bm{s}_{\mathrm{sp}}$
(a natural extension of Definition~\ref{def:dense1d}).

The collection of dense intervals over all spatial start points,
\[
  \mathcal{R}_\theta(P)
  = \{(\bm{s}_{\mathrm{sp}},\, [a_j,\, b_j + W_0])\},
\]
is called the \emph{dense region} of $P$.
\end{definition}

\paragraph{Apriori monotonicity.}

\begin{lemma}[Monotonicity of dense-start sets]\label{lem:apriori}
For any pattern $P$ and its subset $Q \subset P$,
and any spatial start vector $\bm{s}_{\mathrm{sp}}$,
\[
  \mathcal{T}_\theta(P,\bm{s}_{\mathrm{sp}})
  \subseteq
  \mathcal{T}_\theta(Q,\bm{s}_{\mathrm{sp}})
\]
holds.
\end{lemma}

The proof is given in Appendix~\ref{sec:proofs}.

By Lemma~\ref{lem:apriori},
the dense-start set of a size-$k$ pattern $P$
is contained within the dense intervals
of each $(k{-}1)$-subset of $P$.
The proposed method exploits this property
by using \emph{intersection of $(k{-}1)$-subset interval maps}
as candidate intervals (\S\ref{sec:alg_k2}).

\subsection{$d$-Dimensional Dense Region Detection Problem}\label{sec:dense_detection}

The $d$-dimensional dense region detection problem is,
given a $d$-dimensional transaction sequence $\mathcal{D}$,
window widths
$\bm{W} = (W_0, W_1, \ldots, W_{d-1}) \in \mathbb{Z}_{>0}^d$,
and a support threshold $\theta$,
to find the complete set of dense regions
$\mathcal{R}_\theta(P)
= \{(\bm{s}_{\mathrm{sp}}, [a_j,\, b_j + W_0])\}$
for each pattern $P$.
The window along each axis $k$ is a closed interval
$[s_k, s_k + W_k]$ ($W_k + 1$ bins).
